% ============================================================
% Chapter 1: Introduction
% Thesis: Solving the FSRCPSP Using Hybrid ML Approaches
% Author: Souvik Chakraborty
% ============================================================
\chapter{Introduction}
\label{chap:introduction}

The Resource-Constrained Project Scheduling Problem (RCPSP) is an
NP-hard combinatorial optimisation problem
\parencite{Blazewicz1983} and arguably the most studied problem in
project scheduling. Early work focused on exact methods such as
integer programming and branch-and-bound
\parencite{Pritsker1969, Mingozzi1998, Brucker1999}, but these only
scale to small and medium-sized instances. Priority-rule heuristics
\parencite{Klein2000, Kolisch1996} offer a fast alternative,
metaheuristics such as genetic algorithms later became the dominant
paradigm \parencite{Hartmann1998, Kolisch2006a}, and more recently
hybrid methods that combine metaheuristic search with exact MIP
components have been proposed \parencite{Pellerin2020}. For a
comprehensive survey, the reader is referred to
\textcite{Hartmann2010}.

Two extensions of the RCPSP are especially relevant for this thesis.
The first is the Stochastic RCPSP (SRCPSP), which introduces
uncertainty, typically by treating activity durations as random
variables \parencite{Moehring1984}. Ignoring such uncertainty
leads to schedules that look good on paper but break down in
execution \parencite{Hartmann2010}. The SRCPSP has been tackled
with proactive scheduling under chance-constraints
\parencite{Bruni2011, Lamas2016} and, more recently, with
reinforcement learning \parencite{Sallam2021}.

The second extension is the Flexible RCPSP (FRCPSP), which considers a
deterministic total workload for individual activities that can be
distributed over time with varying resource requirements. For each
activity~$i$ and resource type~$k$, the resource allocation~$r_{ikt}$ can
be varied during execution within predefined bounds
$[\underline{r}_{ik},\, \overline{r}_{ik}]$, giving rise to flexible
resource profiles and decision-relevant activity durations that
depend on the allocation. The FRCPSP was first studied by
\textcite{Kolisch2003} in the context of real-world pharmaceutical
research projects and later formalised in mixed-integer programming models
by \textcite{Naber2014}.

The combination of stochastic workloads and flexible resource
profiles, to the best of our knowledge unstudied before
\textcite{MendlRieck2024}, gives rise to the Flexible Stochastic
Resource-Constrained Project Scheduling Problem (FSRCPSP).
\textcite{MendlRieck2024} formulate a chance-constrained MIP,
evaluate it in a performance analysis, and propose a serial
schedule generation scheme (SSGS) as a constructive heuristic.

However, the FSRCPSP is computationally expensive: the SAA model
grows as $O(|V| \cdot |K| \cdot |T| \cdot |S|)$, making exact
solution intractable beyond small instances, and the SSGS cannot
exploit the MIP structure. Decomposition heuristics such as
Relax-and-Fix (R\&F) offer a middle ground by solving a sequence
of smaller subproblems \parencite{helber2010, Tritschler2017},
but a natural strategy of fixing one activity at a time limits
the structural context available per subproblem.

This thesis proposes a different decomposition: unsupervised
machine learning partitions activities into semantically coherent
groups, each fixed simultaneously in a single R\&F subproblem,
giving the solver richer combinatorial context for branching and
bounding. As Chapter~\ref{chap:experiments} shows, the
group-based approach outperforms the baseline in both feasibility
rates and runtime.


% ------------------------------------------------------------
\section{Problem Statement and Goal of This Thesis}
\label{sec:problem_statement}
% ------------------------------------------------------------

The goal of this thesis is to develop and evaluate a group-based
activity-oriented R\&F heuristic in which Ward hierarchical
clustering guides the decomposition, and to compare it against a
baseline that fixes one activity per subproblem
\parencite{MendlRieck2024}. The thesis
aims to answer the following research questions:

\begin{enumerate}
  \item Can unsupervised machine learning (specifically, feature
        extraction and hierarchical clustering) be used to guide the
        decomposition of a Relax-and-Fix heuristic for the FSRCPSP, so that
        semantically similar activities are grouped and fixed
        simultaneously?
  \item Does the group-based R\&F heuristic improve solution quality,
        feasibility rates, and runtime efficiency compared to the baseline
        activity-oriented approach that fixes one activity at a time?
  \item Which activity features (network topology, scheduling urgency,
        or stochastic workload statistics) contribute most to the
        clustering quality and ultimately to the scheduling performance?
  \item Do the proposed scalability extensions (time aggregation and
        per-group scenario reduction) enable the method to handle larger
        problem instances effectively?
\end{enumerate}

The main contributions of this thesis are:
\begin{enumerate}
  \item A group-based R\&F heuristic that uses Ward clustering to
        partition activities into precedence-safe groups, each
        fixed simultaneously in a single MIP subproblem.
  \item A candidate feature pool of $18 + 8|K|$ activity
        descriptors spanning network topology, scheduling urgency,
        and stochastic workload statistics, together with a
        comprehensive ablation study that identifies the most
        informative subset.
  \item Two scalability extensions, time aggregation and per-group
        scenario reduction via PAM K-medoids, that reduce
        per-subproblem model size by up to $10\times$ and enable
        the method to handle instances with $n = 120$ activities.
  \item A systematic evaluation on PSPLIB-based instances with
        $n \in \{10, 30, 60, 120\}$ and $|K| \in \{1, 2, 3, 4\}$,
        demonstrating improvements in feasibility and runtime.
\end{enumerate}


% ------------------------------------------------------------
\section{Structure of This Thesis}
\label{sec:structure}
% ------------------------------------------------------------

The remainder of this thesis is organised as follows.
Chapter~\ref{chap:relatedwork} reviews the deterministic RCPSP,
the SRCPSP, and the FRCPSP, together with the emerging literature
on machine learning for combinatorial optimisation, and positions
our work within that landscape.
Chapter~\ref{chap:formulation} presents the chance-constrained
model that integrates stochastic and flexible characteristics via a
sample average approximation.
Chapter~\ref{chap:baseline} introduces a serial schedule generation
scheme for warm-starting and a baseline activity-oriented
Relax-and-Fix heuristic.
Chapter~\ref{chap:grouprf} extends this baseline with a group-based
decomposition that uses feature extraction and Ward hierarchical
clustering to partition activities into structurally coherent
groups.
Chapter~\ref{chap:experiments} reports computational results on
PSPLIB-based instances with $n \in \{10, 30, 60, 120\}$ activities and
$|K| \in \{1, 2, 3, 4\}$ resource types
\parencite{KolischSprecher1997}, including an ablation study on the
contribution of individual feature groups.
Finally, Chapter~\ref{chap:conclusion} summarises the findings and
discusses directions for future research.
